%----------------------------------------------------------------------------------------
%   USEFUL COMMANDS
%----------------------------------------------------------------------------------------

\newcommand{\dipartimento}{Dipartimento di Matematica ``Tullio Levi-Civita''}

%----------------------------------------------------------------------------------------
% 	USER DATA
%----------------------------------------------------------------------------------------

% Data di approvazione del piano da parte del tutor interno; nel formato GG Mese AAAA
% compilare inserendo al posto di GG 2 cifre per il giorno, e al posto di 
% AAAA 4 cifre per l'anno
\newcommand{\dataApprovazione}{Data}

% Dati dello Studente
\newcommand{\nomeStudente}{Lorenzo}
\newcommand{\cognomeStudente}{Soligo}
\newcommand{\matricolaStudente}{2101057}
\newcommand{\emailStudente}{lorenzo.soligo@studenti.unipd.it}
\newcommand{\telStudente}{+ 39 377 523 1777}

% Dati del Tutor Aziendale
\newcommand{\nomeTutorAziendale}{Gregorio}
\newcommand{\cognomeTutorAziendale}{Piccoli}
\newcommand{\emailTutorAziendale}{gregorio.piccoli@mail.it}
\newcommand{\telTutorAziendale}{+ 39 0371 59 457 24}
\newcommand{\ruoloTutorAziendale}{}

% Dati dell'Azienda
\newcommand{\ragioneSocAzienda}{Zucchetti S.p.A}
\newcommand{\indirizzoAzienda}{Viale della Navigazione Interna 91, Noventa Padovana (PD)}
\newcommand{\sitoAzienda}{http://zucchetti.it/}
\newcommand{\emailAzienda}{info@zucchetti.it}
\newcommand{\partitaIVAAzienda}{P.IVA 12345678999}

% Dati del Tutor Interno (Docente)
\newcommand{\titoloTutorInterno}{Prof.}
\newcommand{\nomeTutorInterno}{Michele}
\newcommand{\cognomeTutorInterno}{Scquizzato}

\newcommand{\prospettoSettimanale}{
     % Personalizzare indicando in lista, i vari task settimana per settimana
     % sostituire a 40 il totale ore della settimana
    \begin{itemize}
        \item \textbf{Prima Settimana - Primo batch di analisi (40 ore)}\\
        Analisi dei seguenti algoritmi di Machine Learning:
        \begin{itemize}
            \item Linear Regression
            \item Logistic Regression
            \item Support Vector Machine
        \end{itemize}
        \item \textbf{Seconda Settimana - Secondo batch di analisi (40 ore)} 
        \\ Analisi dei seguenti algoritmi di Machine Learning:
        \begin{itemize}
            \item Decision Tree
            \item Random Forest
            \item XGBoost
        \end{itemize}
        \item \textbf{Terza Settimana - Prompt Engineering sugli algoritmi (40 ore)} 
        \\ Preparazione dei prompt per LLM per i primi 6 algoritmi analizzati al fine di rendere spiegabili i risultati ottenuti da questi algoritmi.
    
        \item \textbf{Quarta Settimana - Test e documentazione (40 ore)} 
        \\ Elaborazione dei risultati ottenuti e documentazione dei progressi fino a questo punto.

        \item \textbf{Quinta Settimana - Terzo batch di analisi (40 ore)} 
        \\Analisi dei seguenti \textit{ensemble} di algoritmi di Machine Learning:
        \begin{itemize}
            \item RuleFit
            \item Approfondimento sul Random Forest
        \end{itemize}
        \item \textbf{Sesta Settimana - Quarto batch di analisi (40 ore)} 
        \\Analisi dei seguenti \textit{ensemble} di algoritmi di Machine Learning:
        \begin{itemize}
            \item Symbolic Regression
        \end{itemize}
        \item \textbf{Settima Settimana - Prompt Engineering sugli ensemble di algoritmi (40 ore)} 
        \\ Preparazione dei prompt per LLM per i 3 \textit{ensemble} di algoritmi analizzati, con particolare attenzione a come l'unione di eterogenea o omogenea possa influenzare i risultati e la spiegabilità.

        \item \textbf{Ottava Settimana - Test e documentazione (40 ore)} 
        \\ Elaborazione dei risultati ottenuti e documentazione dei progressi fino a questo punto.

    \end{itemize}
}

% Indicare il totale complessivo (deve essere compreso tra le 300 e le 320 ore)
\newcommand{\totaleOre}{320}

\newcommand{\obiettiviObbligatori}{
	 \item \underline{\textit{O01}}: Comprensione profonda degli algoritmi scelti;
	 \item \underline{\textit{O02}}: Comprensione delle tecniche di interpretabilità e spiegabilità dei modelli di Machine Learning;
	 \item \underline{\textit{O03}}: Comprensione delle tecniche di design di algoritmi;
	 \item \underline{\textit{O04}}: Preparazione di prompt per LLM per migliorare la comprensione e l'interpretabilità dei risultati ottenuti dagli algoritmi analizzati.
	 \item \underline{\textit{O05}}: Applicazione principi di Prompt Engineering
}

\newcommand{\obiettiviDesiderabili}{
    \item \underline{\textit{D01}}: Creazione di una metrica per valutare l'efficacia dei prompt elaborati in termini di miglioramento della comprensione e dell'interpretabilità dei risultati ottenuti dagli algoritmi analizzati;
    \item \underline{\textit{D02}}: Automatizzazione dell'elaborazione dei risultati ottenuti dagli algoritmi e passaggio all'LLM.
}

\newcommand{\obiettiviFacoltativi}{
	\item \underline{\textit{F01}}: Applicazione in ambiti reali dei principi di interpretabilità e spiegabilità dei modelli di Machine Learning;
    \item \underline{\textit{F02}}: Studio della spiegabilità in ambito di algoritmi di semi-supervised learning;
}